\documentclass[12pt,a4paper]{article}
\usepackage{fontspec}
\usepackage{geometry}
\usepackage{hyperref}
\usepackage{booktabs}
\usepackage{enumitem}
\usepackage{xcolor}

\geometry{margin=2.4cm}
\IfFontExistsTF{Times New Roman}
  {\setmainfont{Times New Roman}}
  {\setmainfont{DejaVu Serif}}
\IfFontExistsTF{Arial}
  {\setsansfont{Arial}}
  {\setsansfont{DejaVu Sans}}
\hypersetup{colorlinks=true,linkcolor=blue,urlcolor=blue}

\title{Thuyết minh phương pháp hệ thống trả lời trắc nghiệm}
\author{Hackaithon MCQ vLLM Submission}
\date{25/06/2026}

\begin{document}
\maketitle

\section{Mục tiêu}

Hệ thống nhận tập câu hỏi trắc nghiệm chính thức tại \texttt{/code/private\_test.json},
suy luận đáp án và ghi kết quả vào \texttt{/code/submission.csv} và
\texttt{/code/submission\_time.csv}. Runtime cũng ghi bản sao vào \texttt{/output}
để thuận tiện debug. Thiết kế ưu tiên
độ chính xác, khả năng tái lập trong container và thời gian suy luận hợp lý
trên môi trường nhiều GPU.

\section{Kiến trúc tổng quan}

Pipeline gồm năm thành phần chính:

\begin{enumerate}[leftmargin=*]
  \item Bộ chuẩn hóa dữ liệu đầu vào, hỗ trợ JSON và CSV với các cột
  \texttt{qid}, \texttt{question}, và các lựa chọn \texttt{A}, \texttt{B}, ...
  \item Router phân loại câu hỏi theo nhóm: tính toán, đọc hiểu, luật/hành
  chính, nhiều lựa chọn, an toàn và câu hỏi tổng quát.
  \item Các solver luật cứng cho những dạng có công thức hoặc pattern ổn định.
  \item vLLM chạy Qwen3.5-4B với LoRA adapter cho các câu cần mô hình ngôn ngữ.
  \item RAG nội bộ bằng BGE-M3 và vector DB cuối cho nhóm luật/hành chính.
\end{enumerate}

\section{Mô hình và suy luận}

Mô hình nền là Qwen3.5-4B chạy qua vLLM để tận dụng continuous batching,
prefix cache và tensor parallelism. Adapter QLoRA được nạp qua cơ chế LoRA của
vLLM. Các prompt yêu cầu mô hình trả lời đúng một nhãn lựa chọn và dòng cuối có
dạng \texttt{<answer>X</answer>}. Bộ parser hậu xử lý trích nhãn cuối cùng và
kiểm tra nhãn hợp lệ theo số lượng lựa chọn của từng câu.

\section{Router và chiến lược theo route}

Router kết hợp heuristic tiếng Việt, dấu hiệu công thức/số học và embedding
BGE-M3. Các route chính:

\begin{itemize}[leftmargin=*]
  \item \textbf{calculation}: ưu tiên solver công thức và answer repair cho câu
  tính toán.
  \item \textbf{reading\_context}: giữ ngữ cảnh gốc, giảm nhiễu và dùng vLLM.
  \item \textbf{law\_admin}: bổ sung RAG bằng vector DB luật/hành chính.
  \item \textbf{multi\_choice\_many}: dùng prompt chặt hơn cho nhiều nhãn.
  \item \textbf{general}: dùng vLLM trực tiếp với format kiểm soát.
\end{itemize}

\section{Law/Admin RAG}

Phiên bản cuối sử dụng \texttt{rag\_vector\_db\_final}. DB này gồm corpus luật
cơ sở và overlay \texttt{known\_law\_admin\_facts}. Khi route là
\texttt{law\_admin}, hệ thống chỉ truy vấn overlay đã kiểm chứng để tránh nhiễu
từ các shard luật rộng. Mỗi chunk overlay chứa QID, câu hỏi, đáp án đã review và
diễn giải ngắn. Trong kiểm thử nội bộ, top-1 retrieval khớp đúng các câu
law/admin đã biết.

\section{Định dạng container}

Container nhận dữ liệu qua volume \texttt{/code} và ghi file dự đoán vào
\texttt{/code/submission.csv}, \texttt{/code/submission\_time.csv}; đồng thời tạo
bản sao \texttt{/output/pred.csv} để dễ debug. Các model lớn không được đóng gói trực tiếp trong
image; thay vào đó được mount tại runtime:

\begin{center}
\begin{tabular}{ll}
\toprule
Mount & Nội dung \\
\midrule
\texttt{/models} & Qwen3.5-4B local Hugging Face model \\
\texttt{/bge/bge-m3} & BGE-M3 embedding model \\
\texttt{/adapters} & QLoRA adapter \\
\texttt{/rag/rag\_vector\_db\_final} & Vector DB RAG cuối \\
\texttt{/code} & private\_test.json, submission.csv, submission\_time.csv \\
\texttt{/data} & public/private test CSV hoặc JSON \\
\texttt{/output} & pred.csv, bản sao submission.csv/submission\_time.csv \\
\bottomrule
\end{tabular}
\end{center}

Theo thông tin BTC bổ sung, môi trường chấm dùng NVIDIA RTX 5060 Ti 16GB
VRAM và RAM 32GB. Lệnh chạy Docker vì vậy khai báo rõ \texttt{--ipc=host}
để vLLM có shared memory ổn định.

\section{Tái lập}

\begin{verbatim}
docker build -t hackaithon-vllm-submission:latest .

docker run --rm --gpus all --ipc=host \
  -v /path/to/test_data:/data:ro \
  -v /path/to/private_test_dir:/code \
  -v /path/to/qwen_models:/models:ro \
  -v /path/to/bge-m3:/bge/bge-m3:ro \
  -v /path/to/adapter_parent:/adapters:ro \
  -v /path/to/rag_vector_db_final:/rag/rag_vector_db_final:ro \
  -v /path/to/output:/output \
  hackaithon-vllm-submission:latest
\end{verbatim}

\section{Kết quả phát triển}

Snapshot public-test mới nhất đạt 429/463 so với reference nội bộ round 4.
Sai số còn lại tập trung chủ yếu ở nhóm general, multi-choice-many và
calculation; nhóm law/admin chỉ còn một câu khác biệt giữa các vòng review nội
bộ.

\end{document}
